% Generated draft: requires mathematical and source review.
\documentclass[11pt]{article}
\usepackage[margin=25mm]{geometry}
\usepackage{fontspec}
\setmainfont{Latin Modern Roman}
\usepackage{amsmath,amssymb,mathtools,unicode-math}
\setmathfont{Latin Modern Math}
\usepackage{enumitem,needspace,xurl,hyperref}
\setlength{\parindent}{0pt}
\setlength{\parskip}{5pt}
\setlength{\emergencystretch}{3em}
\newcommand{\N}{\mathbb{N}}
\newcommand{\Z}{\mathbb{Z}}
\newcommand{\Q}{\mathbb{Q}}
\newcommand{\R}{\mathbb{R}}
\newcommand{\C}{\mathbb{C}}
\newcommand{\F}{\mathbb{F}}
\newcommand{\A}{\mathbb{A}}
\newcommand{\PP}{\mathbb{P}}
\newcommand{\E}{\mathbb{E}}
\newcommand{\eps}{\varepsilon}
\newcommand{\dd}{\,\mathrm{d}}
\begin{document}
\section*{UnsolvedMath ID: 2640}
\textbf{Source problem code:} catalogue entry
\subsection{Definitions and mathematical statement}
The construction request is to produce a group $G$ in the source's class of Golod groups, a subgroup $H\leq G$, an infinite group $T$ satisfying the source's AT condition, and a surjective group homomorphism
\[
\varphi:H\longrightarrow T.
\]
Here homomorphism means $\varphi(h_1h_2)=\varphi(h_1)\varphi(h_2)$ for all $h_1,h_2\in H$, and surjectivity means that for every $t\in T$ there is $h\in H$ with $\varphi(h)=t$. Equivalently, one seeks a normal subgroup $N\trianglelefteq H\leq G$ with $|H/N|$ infinite and $H/N$ an AT-group in the cited sense; then $T$ is isomorphic to $H/N$ and the quotient map supplies the required epimorphism. The domain need not be the whole Golod group, and the map is not requested to be injective.

The supplied record does not state the defining hypotheses of either specialized class. It refers to Kourovka Problem 9.76 for ``Golod group'' and to A. V. Rozhkov, Mathematical Notes 40, no. 5 (1986), 827--836 for ``AT-group,'' with an additional comparison to Problem 13.55. These are source-dependent terms; no replacement by all Golod--Shafarevich groups, no particular prime, no torsion or finite-generation convention, and no expansion of the abbreviation AT is assumed. Until those precise definitions are supplied from the cited references, the quantified group-homomorphism and quotient requirements above are precise, but the two class predicates remain unresolved and a fully self-contained construction problem cannot be stated.

The source attributes the discussion to A. V. Timofeenko. The registry and its literature triage dated 2026-08-17 label the problem open and report no verified subgroup epimorphism onto an infinite AT-group. That no-result-verified report does not establish nonexistence and is not an independent literature review in this draft.
\subsection{Short English statement}
Construct an infinite AT-group as a quotient of a subgroup of a Golod group. The required map is surjective and the subgroup can be proper. The supplied source references, but does not define, the exact Golod and AT classes, so those definitions require source review.
\subsection{Sources}
\begin{itemize}
\item Kourovka Notebook, New Problems, 21st issue (2026), original Problem 21.131, with discussion attributed to A. V. Timofeenko: \url{https://www.kourovka-notebook.org/}.
\item Kourovka Notebook Problem 9.76, supplied cross-reference for the definition of Golod group, and Problem 13.55, supplied comparison; their text is not included in this record.
\item A. V. Rozhkov, Mathematical Notes 40, no. 5 (1986), 827--836, supplied reference for the AT-group definition; no title or direct URL is supplied.
\item Local registry data/unsolvedmath/problems.json, permanent catalogue entry, source code catalogue entry, ``Kourovka Notebook Problem 21.131''; statement, source-specific definition references, and literature triage dated 2026-08-17. Dataset attribution: Ulam AI and the UnsolvedMath Contributors.
\end{itemize}
\end{document}
